\chapter{ชลศาสตร์ท่อและการคำนวณการใช้น้ำของทุเรียนตามระยะฟีโนโลยี}
\label{chap:ch03}

\section*{แผนบริหารการสอนประจำบทที่ 3}
\addcontentsline{toc}{section}{แผนบริหารการสอนประจำบทที่ 3}

\begin{planbox}{หัวข้อเนื้อหาประจำบท}
\begin{enumerate}[topsep=2pt, itemsep=1pt, partopsep=0pt, parsep=0pt, leftmargin=1.5em]
    \item สรีรวิทยาการใช้น้ำและการคายน้ำของทุเรียนตามระยะฟีโนโลยีในภาคตะวันออก
    \item สมการคำนวณการคายระเหยน้ำอ้างอิงและการประเมินค่าสัมประสิทธิ์พืช
    \item การคำนวณปริมาตรน้ำที่ต้องการต่อต้นและระยะเวลาการให้น้ำที่แม่นยำ
    \item ทฤษฎีไฮดรอลิกส์การไหลในท่อปิดและสมการการสูญเสียแรงดันของเฮเซน-วิลเลียมส์
    \item การคำนวณเฮดรวมของระบบและเกณฑ์การเลือกขนาดปั๊มน้ำพลังงานแสงอาทิตย์
\end{enumerate}
\end{planbox}

\begin{planbox}{วัตถุประสงค์เชิงพฤติกรรม}
\begin{enumerate}[topsep=2pt, itemsep=1pt, partopsep=0pt, parsep=0pt, leftmargin=1.5em]
    \item ระบุค่าสัมประสิทธิ์พืชของทุเรียนในแต่ละระยะการพัฒนาการได้อย่างถูกต้อง
    \item คำนวณปริมาณการใช้น้ำของพืชและอัตราการไหลรวมของแปลงเกษตรได้อย่างแม่นยำ
    \item ประเมินการสูญเสียแรงดันเนื่องจากความเสียดทานในท่อเมนและท่อย่อยได้อย่างถูกต้อง
    \item คำนวณค่าเฮดรวมและจุดทำงานของปั๊มน้ำชลศาสตร์ได้อย่างสมบูรณ์
\end{enumerate}
\end{planbox}

\newpage

\begin{planbox}{วิธีสอนและกิจกรรมการเรียนการสอน}
\begin{enumerate}[topsep=2pt, itemsep=1pt, partopsep=0pt, parsep=0pt, leftmargin=1.5em]
    \item บรรยายหลักการสรีรวิทยาพืชชลศาสตร์และการไหลของของไหลในท่อ
    \item นำเสนอตัวอย่างการคำนวณการใช้น้ำจริงในแปลงวิจัยทุเรียนหมอนทอง จันทบุรี
    \item ปฏิบัติการจำลองขนาดท่อและอัตราแรงดันตกด้วยแบบจำลองคอมพิวเตอร์ Python
    \item กิจกรรมฝึกแก้โจทย์ปัญหาการเลือกขนาดท่อและขนาดปั๊มน้ำให้สอดคล้องกับงบประมาณ
\end{enumerate}
\end{planbox}

\begin{planbox}{สื่อการเรียนการสอน}
\begin{enumerate}[topsep=2pt, itemsep=1pt, partopsep=0pt, parsep=0pt, leftmargin=1.5em]
    \item เอกสารคำสอนบทที่ 3 สไลด์นำเสนอ และตารางค่าคงที่ไฮดรอลิกส์
    \item แบบจำลองสคริปต์ Python คำนวณการสูญเสียแรงดันในท่อ Hazen-Williams
    \item ตัวอย่างท่อและอุปกรณ์ข้อต่อ PE, PVC ขนาด 1 ถึง 3 นิ้ว และหัวจ่ายน้ำมินิสปริงเกลอร์
    \item ซอฟต์แวร์วิเคราะห์เส้นโค้งสมรรถนะปั๊มน้ำ (Pump Performance Curve)
\end{enumerate}
\end{planbox}

\begin{planbox}{การวัดและประเมินผล}
\begin{enumerate}[topsep=2pt, itemsep=1pt, partopsep=0pt, parsep=0pt, leftmargin=1.5em]
    \item การทดสอบการคำนวณปริมาณน้ำตามระยะฟีโนโลยีของทุเรียน
    \item การประเมินความถูกต้องของแบบฝึกหัดการคำนวณการสูญเสียแรงดันในท่อ
    \item การประเมินโครงงานย่อยการออกแบบระบบชลศาสตร์ท่อในแปลงเกษตรตัวอย่าง
\end{enumerate}
\end{planbox}

\clearpage

\section{สรีรวิทยาการใช้น้ำของทุเรียนตามระยะฟีโนโลยี}

การจัดการน้ำในสวนทุเรียนถือเป็นหัวใจสำคัญสูงสุดในการกำหนดผลผลิตและคุณภาพ ทุเรียนเป็นไม้ผลที่มีความไวต่อสภาพการขาดน้ำและน้ำขัง การให้น้ำไม่ถูกต้องในแต่ละระยะการเจริญเติบโตจะส่งผลให้ดอกร่วง ผลอ่อนสลัดทิ้ง เกิดอาการเนื้อแกน หรือทุเรียนสุกไม่สม่ำเสมอ

\subsection{สมการการคายระเหยน้ำและการประเมินค่าสัมประสิทธิ์พืช}

ปริมาณการใช้น้ำของพืชคำนวณจาก \textbf{การคายระเหยน้ำของพืช (Crop Evapotranspiration หรือ $ET_c$)} ตามแนวทางขององค์การอาหารและเกษตรแห่งสหประชาชาติ (FAO-56 Penman-Monteith) ดังสมการ

\begin{equation}
ET_c = ET_0 \times K_c \quad (\text{mm/day})
\label{eq:etc}
\end{equation}

โดยที่ $ET_0$ แทนอัตราการคายระเหยน้ำอ้างอิงของพืชมาตรฐานที่ประเมินจากสถานีอุตุนิยมวิทยา และ $K_c$ แทน \textbf{ค่าสัมประสิทธิ์พืช (Crop Coefficient)} ซึ่งเปลี่ยนแปลงไปตามระยะพัฒนาการทางฟีโนโลยี (Phenological Stages) ของทุเรียน ดังตารางต่อไปนี้

\begin{table}[H]
\centering
\caption{ค่าสัมประสิทธิ์พืช ($K_c$) และปริมาณน้ำแนะนำสำหรับทุเรียนทรงพุ่ม 6 เมตร ในภาคตะวันออก}
\label{tab:durian_kc}
\begin{tabularx}{\textwidth}{p{3.2cm}cccY}
    \toprule
    \textbf{ระยะพัฒนาการ} & \textbf{ช่วงเดือน} & \textbf{ค่า $K_c$} & \textbf{$ET_0$ เฉลี่ย (mm/day)} & \textbf{ปริมาณน้ำแนะนำ (ลิตร/ต้น/วัน)} \\
    \midrule
    1. ฟื้นฟูต้นหลังเก็บเกี่ยว & ส.ค. - ก.ย. & 0.85 & 4.2 & 90 - 110 \\
    2. กักน้ำสร้างตาดอก & ต.ค. - พ.ย. & 0.30 & 4.8 & 0 - 35 (กักน้ำงดชลประทาน) \\
    3. เหยียดตาดอกถึงหางแย้ & ธ.ค. - ม.ค. & 0.70 & 4.5 & 80 - 100 \\
    4. พัฒนาการผลอ่อน (ไข่ไก่) & ก.พ. - มี.ค. & 0.95 & 5.2 & 130 - 160 \\
    5. ผลขยายตัวเต็มที่ & มี.ค. - เม.ย. & 1.10 & 5.8 & 180 - 220 \\
    6. ผลแก่ก่อนเก็บเกี่ยว 15 วัน & พ.ค. & 0.60 & 5.0 & 60 - 80 (ลดน้ำเพื่อสร้างเนื้อ) \\
    \bottomrule
\end{tabularx}
\end{table}

\subsection{การคำนวณปริมาตรน้ำที่ต้องการต่อต้น}

ปริมาตรน้ำที่ต้องให้แก่ต้นทุเรียนต่อวันคำนวณจากพื้นที่ใต้ทรงพุ่มและประสิทธิภาพของระบบให้น้ำ ดังความสัมพันธ์

\begin{equation}
V_{\text{water}} = \frac{A_{\text{canopy}} \times ET_c}{\text{Eff}} \quad (\text{Liters/tree/day})
\label{eq:vwater}
\end{equation}

โดยที่ $A_{\text{canopy}} = \frac{\pi}{4} D_{\text{canopy}}^2$ แทนพื้นที่เงาทรงพุ่มในหน่วยตารางเมตร, $D_{\text{canopy}}$ แทนเส้นผ่านศูนย์กลางทรงพุ่ม, และ $\text{Eff}$ แทนประสิทธิภาพของระบบให้น้ำมินิสปริงเกลอร์ (โดยทั่วไปกำหนดที่ร้อยละ 85 หรือ $0.85$)

\section{ทฤษฎีชลศาสตร์การไหลในท่อและการสูญเสียแรงดัน}

การส่งน้ำจากแหล่งน้ำไปยังหัวสปริงเกลอร์ประจำต้นทุเรียนในแปลงขนาดใหญ่ จำเป็นต้องออกแบบขนาดเส้นผ่านศูนย์กลางท่อให้เหมาะสม เพื่อไม่ให้เกิดการสูญเสียแรงดันเกินพิกัดและรักษาอัตราการไหลให้สม่ำเสมอทุกโซน

\subsection{สมการการสูญเสียแรงดันของเฮเซน-วิลเลียมส์}

การสูญเสียแรงดันเนื่องจากความเสียดทานในท่อปิดภายใต้การไหลแบบปั่นป่วนของน้ำ คำนวณด้วย \textbf{สมการเฮเซน-วิลเลียมส์ (Hazen-Williams Formula)\index{สมการเฮเซน-วิลเลียมส์ (Hazen-Williams)}\index{Hazen-Williams Equation}}

\begin{equation}
h_f = 10.67 \times \frac{L \times Q^{1.852}}{C^{1.852} \times D^{4.87}} \quad (\text{meters})
\label{eq:hazen_williams}
\end{equation}

โดยที่ $h_f$ แทนการสูญเสียแรงดันเนื่องจากความเสียดทาน (Head Loss, m), $L$ แทนความยาวท่อ (m), $Q$ แทนอัตราการไหลในท่อ ($\text{m}^3/\text{s}$), $D$ แทนเส้นผ่านศูนย์กลางภายในท่อ (m), และ $C$ แทนสัมประสิทธิ์ความเรียบของผนังท่อ (สำหรับท่อ PE และ PVC กำหนดค่า $C = 150$)

\begin{theorybox}[เกณฑ์ความเร็วการไหลที่เหมาะสมในระบบท่อชลศาสตร์]
ความเร็วการไหลของน้ำในท่อเมนและท่อย่อยที่เหมาะสมต้องอยู่ระหว่าง \textbf{1.5 ถึง 2.0 เมตรต่อวินาที}:
\begin{itemize}[leftmargin=1.5em]
    \item หากความเร็วต่ำกว่า $0.6\text{ m/s}$ ตะกอนแขวนลอยจะตกสะสมในเส้นท่อ
    \item หากความเร็วสูงกว่า $2.5\text{ m/s}$ การสูญเสียแรงดันจากความเสียดทานจะพุ่งสูงขึ้นอย่างทวีคูณ และเกิดความเสี่ยงรุนแรงต่อการเกิดค้อนน้ำ (Water Hammer) ทำให้ท่อฉีกแตกเสียหาย
\end{itemize}
\end{theorybox}

\subsection{การคำนวณเฮดรวมของระบบเพื่อเลือกปั๊มน้ำ}

\textbf{เฮดรวมของระบบ (Total Dynamic Head หรือ TDH)\index{เฮดรวมของระบบ (TDH)}\index{Total Dynamic Head (TDH)}} ประกอบด้วย 4 องค์ประกอบหลัก ได้แก่

\begin{equation}
TDH = H_s + h_f + h_m + H_{op} \quad (\text{meters})
\label{eq:tdh}
\end{equation}

โดยที่ $H_s$ แทนเฮดสถิต (Static Head: ผลต่างความสูงระหว่างผิวน้ำกับจุดจ่ายน้ำสูงสุด), $h_f$ แทนการสูญเสียแรงดันในท่อเมนและท่อย่อย, $h_m$ แทนการสูญเสียแรงดันรองในข้อต่อและวาล์ว (Minor Loss ประมาณ 10\% ของ $h_f$), และ $H_{op}$ แทนแรงดันใช้งานของหัวจ่ายมินิสปริงเกลอร์ (โดยทั่วไปต้องการ $15 - 20\text{ m}$ หรือ $1.5 - 2.0\text{ bar}$)

\begin{pythonbox}[โปรแกรมจำลองการคำนวณไฮดรอลิกส์และเลือกขนาดท่อน้ำทุเรียน]
import math

def calculate_pipe_hydraulics(flow_rate_lpm: float, pipe_dn_mm: float, length_m: float, c_factor: float = 150.0):
    # แปลงหน่วยเป็น m3/s และ m
    Q = (flow_rate_lpm / 1000.0) / 60.0
    # หนาชั้นท่อ PVC ชั้น 8.5 โดยประมาณ
    inner_diameter_m = (pipe_dn_mm - 4.0) / 1000.0
    area = (math.pi / 4.0) * (inner_diameter_m ** 2)
    velocity = Q / area
    
    # คำนวณ Friction Loss ตามสมการ Hazen-Williams
    hf = 10.67 * length_m * (Q ** 1.852) / ((c_factor ** 1.852) * (inner_diameter_m ** 4.87))
    
    return {
        "velocity_ms": velocity,
        "head_loss_m": hf,
        "is_safe_velocity": 0.8 <= velocity <= 2.2
    }

# ตัวอย่าง คำนวณท่อเมนขนาด 65 มม. ยาว 120 เมตร อัตราไหล 250 ลิตรต่อนาที
# result = calculate_pipe_hydraulics(250.0, 65.0, 120.0)
\end{pythonbox}

\section*{คำถามทบทวนท้ายบทที่ 3}
\addcontentsline{toc}{section}{คำถามทบทวนท้ายบทที่ 3}

\begin{enumerate}
    \item เหตุใดในช่วงพัฒนาการผลขยายตัวเต็มที่ของทุเรียน ค่าสัมประสิทธิ์พืช ($K_c$) จึงสูงถึง 1.10 ในขณะที่ช่วงกักน้ำเพื่อสร้างตาดอกต้องลดค่าลงเหลือ 0.30
    \item สวนทุเรียนแปลงหนึ่งมีขนาดทรงพุ่มเฉลี่ย 7 เมตร อัตราการคายระเหยน้ำ $ET_0 = 5.0\text{ mm/day}$ และ $K_c = 0.95$ จงคำนวณหาปริมาตรน้ำที่ต้องให้ต่อต้นต่อวัน เมื่อระบบสปริงเกลอร์มีประสิทธิภาพ 85\%
    \item หากเพิ่มอัตราการไหลในท่อเดิมขึ้น 2 เท่า การสูญเสียแรงดันเนื่องจากความเสียดทานจะเพิ่มขึ้นประมาณกี่เท่าตามสมการเฮเซน-วิลเลียมส์
    \item จงอธิบายองค์ประกอบของค่าเฮดรวม ($TDH$) และวิเคราะห์ผลกระทบหากเลือกปั๊มน้ำที่มีค่าเฮดต่ำกว่าค่าคำนวณจริง
\end{enumerate}
\clearpage
